\begin{abstract}
This paper examines applications of Jacobian triangularity on invariant
faces corresponding to siphon sets for positive production networks,
encompassing higher-order SIS networks, simplicial contagion models
\cite{IPB}, the threshold models of Granovetter, Kaye and Boudourides
\cite{granovetter1978,Kaye26,Boud26}, and the firing-rate networks of
neuroscience \cite{Betteti2025}. Some main ideas are:
\begin{enumerate}
\ideaitem{1}{First-order production terms in the absent variables determine
the transverse invasion matrices and invasion numbers on siphon faces, while
{\bf terms at least quadratic in the absent variables} leave the transverse
matrix unchanged at a fixed resident state, but can create bistability.\idearest{
Theorem~\ref{t:siphMulti} constructs, for each prescribed finite number,
a smooth response vanishing on $(-\infty,0]$ and positive on $(0,\infty)$
that produces at least that many stable positive equilibria, while
preserving the entire siphon lattice. The response is flat at onset, so production does not contribute to any
transverse Jacobian, which reduces to the linear removal terms.}}

\ideaitem{4}{For threshold networks with a prescribed response-slope
bound $L$, the next-generation matrix method extends to a sharp uniform
test $\bar R_0<1/L$.}
With $\bar R_0=\rho(\beta WV^{-1})$, Theorem~\ref{t:KBsharp} shows that
\[
 \begin{cases}
 \bar R_0<\dfrac1L: & \parbox[t]{.7\linewidth}{every response in the class gives a unique globally exponentially attracting equilibrium;}\\[12pt]
 \bar R_0>\dfrac1L: & \parbox[t]{.7\linewidth}{for each $m\ge2$, some response in the class gives at least $m$ locally exponentially stable equilibria.}
 \end{cases}
\]
The test is uniform over the
response class; the constructive result is that this threshold is sharp
for the possibility of arbitrarily many stable equilibria.

\ifdefined\withSpatial
\ideaitem{spat}{With diffusion, at a spatially homogeneous resident on a siphon face, with irreducible transverse Metzler matrix and nonzero diffusion in that block, $R_\Sigma>1$ (an unstable resident) gives the positive linear \textbf{spreading speed} $c_*=\inf_{\lambda>0}s(\lambda^2\mathcal D+F_\Sigma-V_\Sigma)/\lambda$.\idearest{
Proposition~\ref{t:spread} gives the minimizer, and Remark~\ref{r:movNGM} turns the speed into a family of next-generation
tests: $c>c_*$ exactly when $\rho\bigl(F_\Sigma(c\lambda I+V_\Sigma-\lambda^2\mathcal D)^{-1}\bigr)<1$ for some admissible $\lambda>0$.
With drift there is a second transition, where the perturbation grows but cannot spread upstream.
For threshold responses with slope bound $L$, Theorem~\ref{t:SpdBd} keeps the transition at $\bar R_0=1/L$ under diffusion, with the
\textbf{uniform speed bound} $\bar c(L)=\inf_{\lambda>0}s(\lambda^2\mathcal D+\beta LW-V)/\lambda$; all the constructed counts of stable equilibria survive diffusion (Corollary~\ref{c:DifSta}); numerically, fronts can invade a linearly stable resident, and in the network with non-disjoint siphons the shared species ignites where two residents meet (Figure~\ref{fig:overFront}).}}
\fi

\ideaitem{3}{{\bf Invariant group partitions} give exact reduced equations.\idearest{
For cooperative networks, Theorem~\ref{t:Spe} shows that these reductions
preserve the spectral abscissa and local exponential equilibrium
stability.
With a simple zero eigenvalue and stable complementary directions,
they also lift saddle-node and cusp bifurcations to the full network.
The cusp of a homogeneous branch with a unimodal response density gives
{\bf hysteresis} in the stimulus (Figure~\ref{fig:hyst}).}}

\ideaitem{odds}{{\bf Odds-polynomial responses} $\Phi=\varphi/(1+\varphi)$, zero below
the onset, preserve the siphon lattice as their coefficients vary.\idearest{
At a fixed resident state, only the first coefficient enters the
transverse invasion matrix; higher coefficients can affect invasion
thresholds indirectly by moving the resident equilibrium.
Proposition~\ref{p:Qua}  establishes that, for the scalar equation
$\dot x=-x+\Phi(x)$, degree four
is the exact minimum for coexistence of a stable origin and a positive
nondegenerate cusp, and gives an explicit quartic attaining it.
Theorem~\ref{t:Net} preserves this cusp under multiplicative network
coupling, together with two stable equilibria and a transversely stable
proper siphon face.}}

\ideaitem{overl}{Siphons are read off the zero sets of the responses and the monomials of the
polynomials $P_i$, and minimal siphons may overlap; adding cross-couplings preserves them
if every added production monomial contains at least
one variable from each such siphon containing the species being produced.\idearest{
Example~\ref{e:KBoverlap}, with $P_1=x_2$, $P_2=x_1x_3$, $P_3=x_2$, has the non-disjoint
minimal siphons $\{1,2\}$ and $\{2,3\}$.
Theorem~\ref{t:KBoverlapCusp} constructs a three-population odds network
with two overlapping minimal siphons whose proper faces carry nonzero
resident equilibria with distinct, explicitly computable invasion numbers.
At the constructed parameter, eight hyperbolic stable equilibria coexist
with a nondegenerate interior cusp; Corollary~\ref{c:Nin} gives at least
nine stable equilibria in a nearby open parameter region.
Non-disjoint siphons allow a \textbf{double boundary transcritical relay}: the production term $\beta x_1x_3$ of the
shared species $2$ enters both transverse matrices while both residents stay fixed, so $\beta$ alone
takes the two residents through transcritical bifurcations in turn (Remark~\ref{r:KBrelay}, Figure~\ref{fig:relay}). With disjoint siphons, the production
of one species enters only one transverse matrix.}}

\ideaitem{5}{Positive weights give {\bf extinction certificates} under a
Metzler comparison inequality and characterize uniform exponential
stability of cooperative variational equations over a compact invariant
set.\idearest{
Theorem~\ref{t:EqWgt} gives a compact-domain consequence of converse
contraction theory \cite{KawanoBesselinkCao2020}: for a
cooperative system on a compact invariant box, global convergence to
a Hurwitz equilibrium yields contraction in a state-dependent
positively weighted norm. This includes convergence to extinction.
\textbf{Lyapunov--Volterra functions are weighted sums of Fenchel--Bregman divergences}, with generator $x\log x-x$, and
response-integral energies use the response-dependent convex generator $h$ with $h'=\Phi^{-1}$
(Subsection~\ref{s:FenBre}); combined with Perron-weighted linear terms, they certify
explicit subsets of the basins of attraction of boundary residents
(Propositions~\ref{p:ComBas} and~\ref{p:LVboundary}).}}
\end{enumerate}
\end{abstract} 